% Build from this directory with: latexmk -xelatex report.tex
% Fonts: Liberation Serif, Liberation Sans, Liberation Mono.
\documentclass[12pt,a4paper]{article}

% Shared formatting for an article compiled with XeLaTeX.
% Load after \documentclass[12pt,a4paper]{article}.
% Fonts: Liberation Serif, Liberation Sans, Liberation Mono.

\usepackage{fontspec}
\setmainfont{Liberation Serif}
\setsansfont{Liberation Sans}
\setmonofont[Scale=0.9]{Liberation Mono}
\usepackage{polyglossia}
\setmainlanguage{russian}
\setotherlanguage{english}

\usepackage[a4paper,margin=17mm]{geometry}
\usepackage{amsmath,amssymb,amsthm}
\usepackage{microtype}
\usepackage{booktabs}
\usepackage{algorithm,algorithmic}
\usepackage{tikz}
\usetikzlibrary{arrows.meta,calc,positioning}
\usepackage[section]{placeins}
\usepackage{titlesec}
\usepackage{enumitem}
\usepackage[hidelinks,unicode]{hyperref}

\urlstyle{same}
\floatname{algorithm}{Алгоритм}
\renewcommand{\algorithmicrequire}{\textbf{Input:}}

\setlength{\parindent}{0pt}
\setlength{\parskip}{0.4em}
\linespread{1.02}
\setlength{\emergencystretch}{0.7em}
\setlist[itemize]{topsep=2pt,itemsep=1pt,parsep=0pt,leftmargin=1.2em}
\setlist[enumerate]{topsep=2pt,itemsep=1pt,parsep=0pt,leftmargin=1.5em}
\titleformat{\section}{\large\bfseries}{\thesection.}{0.5em}{}
\titlespacing*{\section}{0pt}{1.0em}{0.3em}
\titleformat{\subsection}{\normalsize\bfseries}{\thesubsection.}{0.5em}{}
\titlespacing*{\subsection}{0pt}{0.7em}{0.15em}
\setcounter{secnumdepth}{2}
\titleformat{\subsubsection}[block]{\normalsize\bfseries}{}{0pt}{}
\titlespacing*{\subsubsection}{0pt}{0.75em}{0.25em}
\pagestyle{plain}
\raggedbottom
\allowdisplaybreaks[2]
\widowpenalty=8000
\clubpenalty=8000
\numberwithin{equation}{section}

\newtheorem{theorem}{Теорема}[section]
\newtheorem{lemma}[theorem]{Лемма}
\newtheorem{corollary}[theorem]{Следствие}
\renewcommand{\proofname}{Доказательство}
\newcommand{\R}{\mathbb R}
\newcommand{\norm}[1]{\lVert #1\rVert}
\newcommand{\ip}[2]{\langle #1,#2\rangle}
\newcommand{\ball}{\overline B}
\newcommand{\dist}{\operatorname{dist}}
\newcommand{\ceil}[1]{\left\lceil #1\right\rceil}
\newcommand{\OO}{\widetilde{\mathcal O}}
\newcommand{\code}[1]{\texttt{\detokenize{#1}}}


\hypersetup{
  pdftitle={Слайдинг при обобщённой гладкости},
  pdfauthor={Исследовательский проект SLIDE}
}

\newcommand{\pass}{\textsf{Pass}}
\newcommand{\classic}{\textsf{RecursiveProcedure}}
\tikzset{
  slide arrow/.style={-{Latex[length=2mm]},thin},
  slide box/.style={draw,align=center,inner sep=4pt,font=\small},
  slide vertex/.style={circle,fill,inner sep=1.6pt}
}

\begin{document}

\begin{center}
  {\Large\bfseries Слайдинг при обобщённой гладкости}\\[2pt]
  {\large Монотонные вариационные неравенства}\\[-1pt]
  {\large и раздельные оценки оракульной сложности}\\[6pt]
  Исследовательский отчёт проекта SLIDE\\[2pt]
  26 ноября 2026 г.
\end{center}

Слайдинг позволяет обращаться к разным частям задачи с разной частотой.
Классическая схема Бородич, Гасникова и Ковалёва использует константы
Липшица градиентов, известные до запуска. В проекте SLIDE рассматривается
более широкий класс, в котором допустимая скорость изменения градиента
зависит от его нормы.

\begin{itemize}
\item \textbf{Путь А.} Адаптивная схема решает монотонное вариационное
неравенство с симметричной $(L_0,L_1)$-гладкостью. Она подбирает кривизну
и рабочий радиус по ходу вычислений. Итоговая теорема, включая число
попыток и бюджеты оракулов по парам, проверена в Lean.
\item \textbf{Локализованный HPE.} Внешний проксимальный шаг даёт
разные оценки для градиента и оператора внутри одной пары.
В общей монотонной задаче ведущие члены по точности имеют порядок
$\varepsilon^{-1/2}$ и $\varepsilon^{-1}$ соответственно, с точностью
до логарифмов. Зависимость от $L_{1,\max}R$ полиномиальная.
\item \textbf{Линейная сходимость.} Если невязка контролирует расстояние
до множества решений, HPE достигает заданного расстояния за логарифмическое
по точности число шагов после учёта конечного бюджета длинных шагов.
Для вычисленной проекционной поправки и критерия остановки есть
формальные доказательства.
\end{itemize}

\setcounter{tocdepth}{1}
\tableofcontents

\clearpage
\section{Постановка и классический бейзлайн}\label{sec:setup}

Пусть $C\subseteq\R^d$ непусто, замкнуто и выпукло, а $x_0\in C$.
У нас есть $n\ge1$ пар $(p_i,Q_i)$. Каждая $p_i\colon\R^d\to\R$
выпукла и дифференцируема. Каждый оператор $Q_i$ монотонен и
$M_i$-липшицев на $C$. Иными словами, для $x,y\in C$
\[
 \ip{Q_i(x)-Q_i(y)}{x-y}\ge0,\qquad
 \norm{Q_i(x)-Q_i(y)}\le M_i\norm{x-y}.
\]
Положим $p=\sum_i p_i$, $Q=\sum_i Q_i$ и $F=\nabla p+Q$.
Ищем решение $z^\star\in C$ вариационного неравенства
\begin{equation}\label{eq:vi}
 \ip{F(z^\star)}{w-z^\star}\ge0\quad\text{для всех }w\in C .
\end{equation}
Предполагаем, что решение существует. Условие означает, что в любой
допустимой точке сравнения направление от $z^\star$ образует с
$F(z^\star)$ угол не больше прямого. При $Q=0$ это условие первого
порядка для минимума $p$ на $C$. В выпукло-вогнутой игре $Q$ может
описывать взаимодействие игроков. Например, билинейный член
$x^\top B y$ даёт оператор $Q(x,y)=(By,-B^\top x)$.

\subsection{Гладкость, проекции и мера ошибки}

Используем симметричную $(L_{0,i},L_{1,i})$-гладкость в смысле
Горбунова и соавторов~\cite{gorbunov}. Для $x,y\in\R^d$
\begin{equation}\label{eq:smooth}
 \norm{\nabla p_i(y)-\nabla p_i(x)}
 \le \bigl(L_{0,i}+L_{1,i}\sup_{\xi\in[x,y]}\norm{\nabla p_i(\xi)}\bigr)
       \norm{y-x}.
\end{equation}
Если $L_{1,i}=0$, это обычная липшицевость градиента. Если $L_{1,i}>0$,
градиент может быть неглобально липшицевым. Например, для $p(t)=\exp(t)$
достаточно $L_0=0$ и $L_1=1$. По теореме о среднем разность
градиентов на отрезке контролируется их максимальной нормой там же.

Вычислительная модель позволяет запрашивать $\nabla p_i$ и $Q_i$
по отдельности. Мы считаем число их вызовов в исходной нумерации.
Значения $p_i$ алгоритм не запрашивает. $\Pi_D$ обозначает точную
евклидову проекцию на непустое замкнутое выпуклое множество $D$.
Она удовлетворяет
$\ip{a-\Pi_D(a)}{w-\Pi_D(a)}\le0$ при $w\in D$.
Проекции на $C$ и на используемые пересечения $C$ с шарами
считаются доступными оракулами. Их стоимость учитывается отдельно.

Качество допустимой точки $z$ измеряем составным зазором
\begin{equation}\label{eq:gap}
 G(z,w)=p(z)-p(w)+\ip{Q(w)}{z-w}.
\end{equation}
Из выпуклости и \eqref{eq:vi} следует $G(z,z^\star)\ge0$ при
$z\in C$. Из монотонности дополнительно следует
$G(z^\star,w)\le0$ при $w\in C$.
Для относительной точности $\rho>0$ требуем
\begin{equation}\label{eq:relative}
 G(z_{\rm out},w)\le\rho\norm{w-x_0}^2\qquad(w\in C).
\end{equation}
Такое условие осмысленно и на неограниченном $C$. На шаре сравнения
$C\cap\ball_R(x_0)$ оно даёт абсолютную точность $\rho R^2$.
При $Q=0$ сравнение с $z^\star$ даёт обычную ошибку по значению
целевой функции.

\subsection{Что переносим из классической статьи}

Основной бейзлайн для общей монотонной задачи --- теорема~4.5 и
следствие~4.6 работы Бородич, Гасникова и Ковалёва~\cite{borodich}.
При известных классических константах липшицевости $L_i$ и $M_i$
они дают относительную оценку зазора и покомпонентный порядок
\[
 \mathcal O\!\left(\max\left\{1,\sqrt{L_i/\rho},\,M_i/\rho\right\}\right)
\]
для обоих оракулов пары при фиксированном $n$. В этой записи
поглощены множители от числа уровней и нормировки точности.
Устройство слайдинга удобно представить на двух парах. Одна итерация
верхнего уровня строит модель первой пары и сохраняет её на весь
внутренний цикл второй пары. При длинах циклов $T_1,T_2$ число
итераций верхнего и внутреннего уровней равно $T_1$ и $T_1T_2$.
Поэтому один оракул можно вызывать существенно реже другого.

Линейная сходимость в той же статье относится к структурированным
билинейным седловым задачам с дополнительными предпосылками
(теорема~4.7). В общей постановке \eqref{eq:vi} линейная скорость
появится в разделе~\ref{sec:linear} после введения отдельного
условия на невязку.

\subsection{Исходный алгоритм слайдинга}\label{sec:classical-algorithm}

В алгоритме~\ref{alg:original} переписан \emph{Algorithm 1} из
приложения~B статьи~\cite{borodich}. Матрица $P\succ0$ задаёт
$\norm{z}_P^2=z^\top Pz$. Вместо многоиндексных функций сохраняем
список $\mathcal M$ уже построенных моделей и аффинное преобразование
$z\mapsto\beta z+b$ оставшихся функций. На входе уровня $k$ такая
функция имеет вид $p_i(\beta z+b)/\beta$ для $i\ge k$.
Вектор $U$ хранит текущие центры и передаётся по ссылке, как в
строках 28--30 оригинала. Запись $[\mathcal M,m]$ добавляет модель
в конец списка. Длины $T_k$ и константы $L_i,M_i$ заданы до запуска.

\begin{algorithm}[H]
\caption{Классический слайдинг Бородич, Гасникова и Ковалёва
в обозначениях отчёта}\label{alg:original}
\begin{algorithmic}[1]
\REQUIRE $x_0\in C$, $P\succ0$, weights $(\alpha_t)$,
  loop lengths $(T_k)$, constants $(L_i,M_i)$ in the $P$-metric.
\STATE $U_i\gets x_0$ for all $i$,
  $s_i\gets\prod_{j=1}^{i}\alpha_{T_j-1}$ for all $i$.
\RETURN $\classic(1,1,0,\varnothing;U)$.
\STATE \textbf{procedure} $\classic(k,\beta,b,\mathcal M;U)$
\IF{$k=n+1$}
  \RETURN $\displaystyle\arg\min_{z\in C}\sum_{m\in\mathcal M}m(z)$.
\ENDIF
\STATE $y\gets U_k$.
\FOR{$t=0,\ldots,T_k-1$}
  \STATE $u\gets U_k$, $\alpha\gets\alpha_t$,
    $\beta'\gets\beta\alpha$, $b'\gets b+\beta(1-\alpha)y$.
  \STATE $g\gets\nabla p_k(\beta'u+b')$, $q\gets Q_k(u)$,
    $H\gets\beta'(L_k+M_k/s_k)$.
  \STATE $m(z)\gets (H/2)\norm{z-u}_P^2+\ip{g+q}{z}$.
  \STATE $v\gets\classic(k+1,\beta',b',[\mathcal M,m];U)$.
  \STATE $y\gets(1-\alpha)y+\alpha v$.
  \STATE $U_k\gets v+(HP)^{-1}(q-Q_k(v))$.
\ENDFOR
\RETURN $y$.
\end{algorithmic}
\end{algorithm}

На рис.~\ref{fig:recursion} показано дерево для $n=3$ и $T_k=2$.
Последняя строка алгоритма возвращает среднее текущего вызова.
Минимизация суммы моделей выполняется в листьях на глубине $n+1$.

\begin{figure}[htbp]
\centering
\begin{tikzpicture}[x=1cm,y=1cm]
  \foreach \a/\b/\c/\d in {
    0/0/-3.2/-1,0/0/3.2/-1,
    -3.2/-1/-4.8/-2,-3.2/-1/-1.6/-2,
    3.2/-1/1.6/-2,3.2/-1/4.8/-2,
    -4.8/-2/-5.6/-3,-4.8/-2/-4/-3,
    -1.6/-2/-2.4/-3,-1.6/-2/-.8/-3,
    1.6/-2/.8/-3,1.6/-2/2.4/-3,
    4.8/-2/4/-3,4.8/-2/5.6/-3}{
    \draw[thin] (\a,\b) -- (\c,\d);
  }
  \node[slide vertex] at (0,0) {};
  \foreach \x in {-3.2,3.2}{\node[slide vertex] at (\x,-1) {};}
  \foreach \x in {-4.8,-1.6,1.6,4.8}{\node[slide vertex] at (\x,-2) {};}
  \foreach \x in {-5.6,-4,-2.4,-.8,.8,2.4,4,5.6}{
    \node[rectangle,draw,fill=white,minimum size=3.2pt,inner sep=0pt] at (\x,-3) {};
  }
  \node[font=\small,anchor=east] at (-6,0) {$k=1$};
  \node[font=\small,anchor=east] at (-6,-1) {$k=2$};
  \node[font=\small,anchor=east] at (-6,-2) {$k=3$};
  \node[font=\small,anchor=east] at (-6,-3) {$k=4$};
  \node[font=\small,anchor=west] at (6,0) {1 вызов};
  \node[font=\small,anchor=west] at (6,-1) {2 вызова};
  \node[font=\small,anchor=west] at (6,-2) {4 вызова};
  \node[font=\small,anchor=west] at (6,-3) {8 минимумов};
\end{tikzpicture}
\caption{Дерево вызовов при $n=3$ и $T_1=T_2=T_3=2$.
Кружок обозначает вызов уровня с собственным циклом,
квадрат обозначает лист.}\label{fig:recursion}
\end{figure}

\clearpage
\section{Путь А. Адаптивная кривизна и рабочий радиус}\label{sec:patha}

Путь А сохраняет рекурсию классического слайдинга. Он выбирает
кривизну модели по тому отрезку, который реально возник на пробном
шаге. При неудачной проверке меняется оценка конкретной пары.
Длины циклов и порядок пар пересчитываются перед следующей попыткой.

\subsection{Алгоритм}

Пусть $s=x_0$, $\delta=\rho/3$, $R_0>0$ и
$\lambda_i=\max\{\delta,\Lambda_i^0\}$ для начальных оценок
$\Lambda_i^0\ge0$. Начинаем с радиуса $R_0$ и кривизн
$\Lambda_i=\lambda_i$. На каждой попытке работаем на
$D=C\cap\ball_R(s)$. Сортируем пары по возрастанию
$h_i=\max\{1,\sqrt{\Lambda_i/\delta},M_i/\delta\}$. В формулах прохода
индексы относятся к этому порядку. Вводим
\begin{align}
 \alpha_0&=1,& \alpha_{t+1}&=\frac{2}{1+\sqrt{1+4/\alpha_t^2}},\\*
 T_1&=\ceil{2h_1},&
 T_i&=\ceil{16h_i/h_{i-1}}\qquad(i>1),\\*
 a_i&=\alpha_{T_i-1},&
 s_i&=\prod_{j=1}^{i}a_j,\qquad c_i=\Lambda_i+2M_i/s_i.
 \label{eq:pa-parameters}
\end{align}
Для списка моделей $m_j(z)=H_j\norm{z-u_j}^2/2+\ip{d_j}{z}$
лист возвращает $\Pi_D\bigl(\sum_j(H_j u_j-d_j)/\sum_j H_j\bigr)$.
Путь А использует евклидову метрику. В алгоритме~\ref{alg:path-a}
глубокие центры хранятся в $U$, как в алгоритме~\ref{alg:original}.
По сравнению с исходным проходом меняются расписание, кривизна
моделей и поправка центра: она использует $\Pi_D$. Проверки секущих
и радиуса могут вызвать повтор всего прохода.

\begin{algorithm}[htbp]
\caption{Путь А. Поиск кривизны и радиуса с проверяемым проходом}
\label{alg:path-a}
\begin{algorithmic}[1]
\REQUIRE $s\in C$, $\rho>0$, $R_0>0$, $(\Lambda_i^0)$, $(M_i)$.
\STATE $R\gets R_0$, $\delta\gets\rho/3$,
       $\Lambda_i\gets\max\{\delta,\Lambda_i^0\}$ for all $i$.
\WHILE{true}
  \STATE $D\gets C\cap\ball_R(s)$.
     Sort the pairs and compute $T_i,s_i,c_i$ using \eqref{eq:pa-parameters}.
  \STATE $U_i\gets s$ for all $i$.
     $y\gets\pass(1,1,0,\varnothing;U)$.
  \IF{the pass rejects pair $i$ with threshold $\tau_i$}
    \STATE $\Lambda_i\gets\max\{2\Lambda_i,\tau_i\}$ using the original pair index.
  \ELSE
    \IF{$\norm{y-s}\le R/2$}
      \RETURN $y$.
    \ENDIF
    \STATE $R\gets2R$.
  \ENDIF
\ENDWHILE
\STATE \textbf{procedure} $\pass(k,\beta,b,\mathcal M;U)$
\IF{$k=n+1$}
  \RETURN $\displaystyle\arg\min_{z\in D}\sum_{m\in\mathcal M}m(z)$ using the projection formula above.
\ENDIF
\STATE $y\gets U_k$.
\FOR{$t=0,\ldots,T_k-1$}
  \STATE $u\gets U_k$, $\alpha\gets\alpha_t$, $\beta'\gets\beta\alpha$,
     $b'\gets b+\beta(1-\alpha)y$.
  \STATE $x\gets\beta'u+b'$, $g\gets\nabla p_k(x)$,
     $q\gets Q_k(u)$, $H\gets\beta'c_k$.
  \STATE $m(z)\gets(H/2)\norm{z-u}^2+\ip{g+q}{z}$.
  \STATE $v\gets\pass(k+1,\beta',b',[\mathcal M,m];U)$.
     Propagate any rejection immediately.
  \STATE $x'\gets\beta'v+b'$.
     Compute $\tau_k$ using \eqref{eq:secant-test}.
  \IF{$\tau_k>\Lambda_k$}
    \RETURN a rejection for pair $k$ with threshold $\tau_k$.
  \ENDIF
  \STATE $y\gets(1-\alpha)y+\alpha v$,
    $U_k\gets\Pi_D\bigl(v+(q-Q_k(v))/H\bigr)$.
\ENDFOR
\RETURN $y$.
\end{algorithmic}
\end{algorithm}

Секущая в проверке вычисляется по формуле
\begin{equation}\label{eq:secant-test}
 \tau_k=\frac{2\ip{\nabla p_k(x')-\nabla p_k(x)}{x'-x}}{\norm{x'-x}^2},
 \qquad \tau_k=0\quad\text{при }x'=x.
\end{equation}

\subsection{Гарантия и цена адаптации}

Пусть $r_\star=\norm{s-z^\star}$,
$\bar R=\max\{R_0,24r_\star\}$ и
\[
 A_i^{\rm ball}(R)=
 \max_{\xi\in C\cap\ball_R(s)}(L_{0,i}+L_{1,i}\norm{\nabla p_i(\xi)}),
 \qquad
 \bar\Lambda_i=\max\{\lambda_i,4A_i^{\rm ball}(\bar R)\}.
\]
Эти величины используются в анализе. Алгоритму достаточно стартового
радиуса, желаемой точности, $M_i$ и доступа к проекциям.

\begin{theorem}[Путь А]\label{thm:patha}
Алгоритм завершается с $z_{\rm out}\in C$ и гарантией
\eqref{eq:relative}. Все попытки имеют радиус не больше $\bar R$
и кривизны $\Lambda_i\le\bar\Lambda_i$. Число попыток $N$, включая
успешную, удовлетворяет
\begin{equation}\label{eq:pa-restart}
 N\le 1+\left\lfloor\log_2(\bar R/R_0)\right\rfloor
       +\sum_i\left\lfloor\log_2(\bar\Lambda_i/\lambda_i)\right\rfloor .
\end{equation}
В исходной нумерации для обоих оракулов пары $i$ имеем
\begin{equation}\label{eq:pa-cost}
 \max\{N_{\nabla p_i},N_{Q_i}\}\le
 6\cdot17^{n-1}N
 \max\left\{1,\sqrt{3\Lambda_i^{\rm final}/\rho},\,3M_i/\rho\right\}.
\end{equation}
\end{theorem}

При $L_{1,i}=0$ можно взять $A_i^{\rm ball}(R)=L_{0,i}$.
С точностью до констант по глубине и множителя числа попыток
восстанавливается классическая зависимость по паре. При
$L_{1,i}>0$ наблюдаемая $\Lambda_i^{\rm final}$ может быть умеренной,
но верхняя оценка $A_i^{\rm ball}(\bar R)$ в общем случае может расти
экспоненциально с $L_{1,i}\bar R$. В формуле \eqref{eq:pa-cost}
оператор и градиент одной пары используют общий максимум.

\subsection{Почему алгоритм работает}

\subsubsection{Сертификат одного шага}

Выпуклость ограничивает остаток линеаризации $p_k$ сверху величиной
\[
 \ip{\nabla p_k(x')-\nabla p_k(x)}{x'-x}.
\]
Поэтому после проверки
$(\tau_k/2)\norm{x'-x}^2$ покрывает этот остаток. В координатах
уровня $f_k(z)=p_k(\beta'z+b')/\beta'$ достаточно кривизны
$L=\beta'\tau_k$. Из $\beta'\ge s_k$ следует $H\ge L+2M_k$.
Для сравнения с $w\in D$ пишем $u^+=\Pi_D(v+(Q_k(u)-Q_k(v))/H)$.
Выпуклость, монотонность и выделение квадрата дают
\begin{align}
 m_k(v)-m_k(w)+\frac H2\norm{v-w}^2
 &\ge f_k(v)-f_k(w)+\ip{Q_k(w)}{v-w}\notag\\
 &\quad+\frac H2\bigl(\norm{u^+-w}^2-\norm{u-w}^2\bigr).
 \label{eq:pa-model}
\end{align}
При выделении квадрата теряется не больше
$M_k^2\norm{v-u}^2/(2H)$. Этот член покрывает запас $H-L$.

\subsubsection{Сложение по уровням}

Зазор смешанной задачи на уровне $k$ имеет вид
\begin{align*}
 \mathcal G_k(z,w)
 &=\sum_{j<k}\left(m_j(z)-m_j(w)+\frac{H_j}{2}\norm{z-w}^2\right)\\
 &\quad+\sum_{i\ge k}\left(
 \frac{p_i(\beta z+b)-p_i(\beta w+b)}{\beta}
 +\ip{Q_i(w)}{z-w}\right).
\end{align*}
В листе он неположителен по условию оптимальности минимума суммы
моделей. На шаге $y^+=(1-\alpha)y+\alpha v$ из выпуклости
и \eqref{eq:pa-model} следует
\[
 \frac{\mathcal G_k(y^+,w)-(1-\alpha)\mathcal G_k(y,w)}{\alpha}
 +\frac H2\bigl(\norm{u^+-w}^2-\norm{u-w}^2\bigr)
 \le\mathcal G_{k+1}(v,w).
\]
Подстановка оценки дочернего уровня и деление ещё на $\alpha$
показывают, что каждому центру соответствует разность квадратов
расстояний. Её коэффициент $(\beta c_i/2)\prod_{j=k+1}^{i}a_j^2$
постоянен вдоль цикла.
Равенство $(1-\alpha_{t+1})/\alpha_{t+1}^2=1/\alpha_t^2$
сокращает соседние зазоры. Умножение на $a_k^2$ после суммирования
даёт оценку для всего вызова уровня $k$. В корне получаем
\begin{equation}\label{eq:pa-energy}
 G(y,w)\le \frac12\sum_i \ell_i
     \bigl(\norm{s-w}^2-\norm{U_i^{\rm out}-w}^2\bigr),
 \qquad \ell_i=s_i^2c_i.
\end{equation}
В частности, эта оценка действует на любом фиксированном выпуклом
рабочем множестве после принятого прохода без проверки выходного радиуса.

\subsubsection{Расписание и цена прохода}

Обратные веса растут за шаг на величину от $1/2$ до $1$,
поэтому $1/T\le\alpha_{T-1}\le2/T$.
Из округления $\prod_{j\le i}T_j\le 3\cdot17^{i-1}h_i$.
Если $a_1h_1=b_1$, то $1/3\le b_1\le1$ и
$s_i h_i\le b_1 8^{-(i-1)}$. Это даёт
$\ell_1\in[\delta b_1^2,3\delta]$ и
$\ell_i\le7\cdot8^{-(i-1)}\ell_1$ для $i>1$.
Значит, $\sum_{i>1}\ell_i\le\ell_1$ и $\frac12\sum_i\ell_i\le3\delta$.
Каждый шаг требует не больше двух вызовов каждого оракула своей пары.
Лист требует одной проекции, как и каждая операторная поправка.
Эти оценки действуют и для оборванной попытки. Суммирование
даёт \eqref{eq:pa-cost}.

\subsubsection{Конечность поиска и глобальный зазор}

Если $z^\star\in D$, то подстановка $w=z^\star$ в корневую оценку
для префикса даёт $\norm{U_1^t-z^\star}\le\sqrt2 r_\star$.
Это следует из неотрицательности зазора и
$\sum_{i>1}\ell_i\le\ell_1$. Из нерастяжимости проекции
$\norm{u^+-v}\le\norm{v-u}/2$, откуда
$\norm{v-z^\star}\le2\norm{u^+-z^\star}+\norm{u-z^\star}$.
Следовательно, выход проходит проверку половины радиуса при
$R\ge12r_\star$. Вместе с удвоением это даёт границу $\bar R$.
На любом пробном отрезке $\tau_k\le 2A_i^{\rm ball}(\bar R)$
для соответствующей исходной пары. Обновляемая координата растёт
хотя бы вдвое и остаётся под своим пределом. Отсюда
\eqref{eq:pa-restart}.

На рабочем шаре \eqref{eq:relative} следует из \eqref{eq:pa-energy}.
Для произвольного $w\in C$ с $d=\norm{w-s}>R$ возьмём
$w'=(1-\theta)y+\theta w$ с $\theta=R/(2d-R)$.
Поскольку $\norm{y-s}\le R/2$, имеем $w'\in D$.
Выпуклость $p$ и монотонность $Q$ дают
$\theta G(y,w)\le G(y,w')$. Тогда
$G(y,w)\le\rho R^2/\theta=\rho R(2d-R)\le\rho d^2$.

\clearpage
\section{Локализованный HPE. Разделение внутри пары}\label{sec:hpe}

Путь А ограничивает кривизну на рабочем шаре, но радиус этого шара
может быть сравним с расстоянием до решения. В HPE
(\emph{hybrid proximal extragradient}) мы делаем серию внешних
проксимальных шагов. Каждый из них решает более простую задачу
на малом шаре, который переносится вместе с текущим центром.
Внешняя схема использует вычисленную невязку подзадачи и потому
допускает неточное внутреннее решение~\cite{monteiro}.

В этом разделе алгоритму известны $L_{0,i},L_{1,i}$ и число $R>0$
с $\norm{x_0-z^\star}\le R$. Положим
\[
 L_{1,\max}=\max_i L_{1,i},\qquad
 r=\min\{R,1/L_{1,\max}\},\qquad
 \mathcal K=1+R^2/r^2,
\]
где $1/0=+\infty$. Рабочее множество на внешнем шаге из $x\in C$
равно $D=C\cap\ball_r(x)$, а регуляризованный оператор равен
\[
 F_\Gamma(z)=F(z)+\Gamma(z-x),\qquad \Gamma>0.
\]
У подзадачи на $D$ есть единственное решение $z_\Gamma^D$.
Существование следует из непрерывности и теоремы Брауэра для
отображения $z\mapsto \Pi_D(z-F_\Gamma(z))$ на компактном $D$.
Единственность следует из $\Gamma$-сильной монотонности.

\begin{figure}[htbp]
\centering
\begin{tikzpicture}[node distance=8mm]
  \node[slide box,text width=31mm] (center) {Внешний центр $x$\\
  $D=C\cap\ball_r(x)$};
  \node[slide box,text width=56mm,right=of center] (probe)
    {Проба с параметром $\Gamma$\\
     Проходы пути А на фиксированном $D$};
  \node[slide box,text width=33mm,right=of probe] (out)
    {Поправка $(u,v)$\\
    Новый центр $x^+$};
  \draw[slide arrow] (center) -- (probe);
  \draw[slide arrow] (probe) -- (out);
  \draw[slide arrow] ([xshift=14mm]probe.south) -- ++(0,-8mm)
    -| node[pos=.25,below,font=\small] {отказ: $\Gamma$ уменьшается}
    ([xshift=-14mm]probe.south);
\end{tikzpicture}
\caption{Один внешний шаг HPE. При отказе меняется регуляризация,
а центр и рабочее множество сохраняются.}\label{fig:hpe-composition}
\end{figure}

\subsection{Внутренний проход и сертификат невязки}

Обозначим $g_i=\nabla p_i$. По \eqref{eq:smooth} коэффициент
$L_{0,i}+L_{1,i}\norm{g_i(\cdot)}$ вдоль отрезка длины $t$
может вырасти не более чем в $\exp(L_{1,i}t)$ раз.
Действительно, градиент липшицев на компактном отрезке,
а разностное неравенство даёт для коэффициента верхнюю оценку
производной $L_{1,i}$, умноженную на текущий коэффициент, при
параметризации по длине. Интегрирование даёт указанную оценку.
Поэтому допустимы вычислимые константы
\[
 \widehat L_i(x)=e\bigl(L_{0,i}+L_{1,i}\norm{g_i(x)}\bigr),\qquad
 L_\Gamma=\Gamma+\sum_i(\widehat L_i(x)+M_i)
\]
для липшицевости $g_i$ и $F_\Gamma$ на $D$, соответственно.

В подзадаче заменяем каждую пару двумя уровнями $(p_i,0)$ и $(0,Q_i)$.
Добавку $\Gamma\norm{z-x}^2/2$ относим к одной гладкой функции.
Её градиент вычисляется напрямую. Применяем фиксированный проход
пути А с $\delta=\Gamma/24$. Для уровня с градиентом, липшицевым
с константой $L$, выбираем кривизну $\max\{\delta,2L\}$. На
операторном уровне берём $L=0$.
Этого достаточно, чтобы ни одна секущая не вызвала отказ.
По \eqref{eq:pa-energy} выход $y$ из начальной точки $a$ даёт
\[
 G_\Gamma(y,w)\le(\Gamma/8)\norm{a-w}^2,\qquad
 G_\Gamma(y,w)=G(y,w)+\frac{\Gamma}{2}
                    \bigl(\norm{y-x}^2-\norm{w-x}^2\bigr).
\]
В частности, $G_\Gamma(y,z_\Gamma^D)\ge\Gamma\norm{y-z_\Gamma^D}^2/2$
по выпуклости исходных функций и условию ВН для $F_\Gamma$.
Значит, каждый такой проход сокращает расстояние до $z_\Gamma^D$
вдвое. При старте из $x$ после
$m_\Gamma=\ceil{\log_2(128L_\Gamma/\Gamma)}$ проходов получаем
\begin{equation}\label{eq:hpe-inner}
 \norm{y-z_\Gamma^D}
 \le \eta_\Gamma:=\frac{\Gamma \norm{z_\Gamma^D-x}}{128L_\Gamma}
 \le \frac r{128}.
\end{equation}
Расстояние $\norm{z_\Gamma^D-x}$ входит только в анализ.
Сохранение этого множителя существенно, когда проксимальное
перемещение мало.

Для принятой пробы вычислим
\[
 u=\Pi_D\bigl(y-F_\Gamma(y)/L_\Gamma\bigr),\qquad
 n_D=L_\Gamma(y-u)-F_\Gamma(y),\qquad
 v=F(u)+n_D .
\]
По свойству проекции $n_D\in N_D(u)$, где
$N_D(u)=\{n:\ip{n}{w-u}\le0,\ w\in D\}$ --- нормальный конус.
Точка $z_\Gamma^D$ неподвижна для того же проекционного отображения.
Нерастяжимость проекции и липшицевость дают
\[
 \norm{u-z_\Gamma^D}\le2\eta_\Gamma,\qquad
 \norm{u-y}\le3\eta_\Gamma,\qquad
 \norm{F_\Gamma(u)+n_D}\le6L_\Gamma\eta_\Gamma .
\]
Если $\norm{u-x}<r$, то $n_D\in N_C(u)$. Для $w\in C$
достаточно подставить в нормальное неравенство
$u+t(w-u)\in D$ при малом $t>0$ и поделить на $t$.
Таким образом, $v\in F(u)+N_C(u)$.
Кроме того, $\norm{u-z_\Gamma^D}\le\norm{z_\Gamma^D-x}/64$,
так что $\norm{u-x}\ge63\norm{z_\Gamma^D-x}/64$.
Подстановка в оценку $\norm{F_\Gamma(u)+n_D}$ даёт
\begin{equation}\label{eq:hpe-res}
 \norm{u-x+\Gamma^{-1}v}\le \frac1{21}\norm{u-x}.
\end{equation}
При $z_\Gamma^D=x$ обе стороны равны нулю.

\subsection{Выбор регуляризации}

В начале внешнего шага берём
$\Gamma^{(0)}=\max\{\Gamma_{\min},8\norm{F(x)}/r\}$.
Если $F(x)=0$, можно вернуть $x$ как решение. На каждой пробе
вычисляем $y$ по \eqref{eq:hpe-inner}. Принимаем первую пробу
с $\norm{y-x}\ge r/4$ или $\Gamma=\Gamma_{\min}$.
При отклонении переходим к следующей пробе с
$\Gamma\gets\max\{\Gamma_{\min},\Gamma/2\}$.
Поправку $u$ вычисляем только для принятой пробы.

Для первой пробы ВН в $z_\Gamma^D$ с точкой сравнения $x$
и монотонность дают
$\Gamma\norm{z_\Gamma^D-x}\le\norm{F(x)}$, если перемещение
ненулевое. Значит, оно не больше $r/8$. При последующих пробах
важно, что $D$ остаётся тем же. Пусть $z_j$ --- точные решения
двух подзадач с $\Gamma_1\ge\Gamma_2>0$ и $d_j=\norm{z_j-x}$.
Если сложить их ВН и применить монотонность и неравенство
Коши--Буняковского, получится
\[
 0\le(d_2-d_1)(\Gamma_1d_1-\Gamma_2d_2).
\]
При $\Gamma_1>\Gamma_2$ это даёт
$d_1\le d_2\le(\Gamma_1/\Gamma_2)d_1$. При равных $\Gamma_j$
решения совпадают. Если предыдущая проба отклонена, её
точное перемещение меньше $r/4+r/128=33r/128$. Следующая
проба увеличивает его максимум вдвое. Поэтому на принятой
пробе $\norm{z_\Gamma^D-x}\le33r/64$, и обе точки $z_\Gamma^D,u$
строго внутри шара. Значит, точное решение регуляризованной
задачи является решением и на $C$, а поправка удовлетворяет
\eqref{eq:hpe-res} с нормалью к $C$. На пробе, принятой по длине,
дополнительно $\norm{u-x}\ge29r/128$.

На рис.~\ref{fig:hpe-radius} показан запас до границы для принятой
пробы. Алгоритм~\ref{alg:hpe} задаёт полный порядок вычислений.
Внутри он вызывает только процедуру \pass{} из алгоритма~\ref{alg:path-a}
на фиксированном $D$. Поиск радиуса пути А здесь не используется.

\begin{figure}[htbp]
\centering
\begin{tikzpicture}[x=12cm,y=1cm]
  \draw[thin] (0,0) -- (1,0);
  \foreach \pos in {0,.25,1}{
    \draw[thin] (\pos,-.09) -- (\pos,.09);
  }
  \node[below,font=\small] at (0,-.09) {$0$};
  \node[below,font=\small] at (.25,-.09) {$r/4$};
  \node[below,font=\small] at (1,-.09) {$r$};
  \draw[thick] (0,.53) -- ({33/64},.53);
  \draw[thin] ({33/64},.45) -- ({33/64},.61);
  \node[font=\small,above] at ({33/128},.53)
    {$\norm{z_\Gamma^D-x}\le33r/64$};
  \draw[thick] (0,1.08) -- ({17/32},1.08);
  \draw[thin] ({17/32},1.00) -- ({17/32},1.16);
  \node[font=\small,above] at ({17/64},1.08)
    {$\norm{u-x}\le17r/32$};
  \draw[dashed,thin] (1,.10) -- (1,1.12);
  \node[font=\small,align=center] at (1,1.39) {граница\\шара};
  \node[font=\small,below] at (.25,-.42) {порог для $\norm{y-x}$};
\end{tikzpicture}
\caption{Радиальные границы на принятой пробе. Отрезки показывают
верхние оценки расстояния до $x$.}\label{fig:hpe-radius}
\end{figure}

\begin{algorithm}[htbp]
\caption{Локализованный HPE с внутренним проходом пути А}
\label{alg:hpe}
\begin{algorithmic}[1]
\REQUIRE $x_0\in C$, $\rho>0$, distance bound $R$, smoothness parameters and $(M_i)$.
\STATE $x\gets x_0$, $S\gets0$, $W\gets0$,
  $\Gamma_{\min}\gets\rho$, $r\gets\min\{R,1/L_{1,\max}\}$.
\WHILE{$S<1/(2\rho)$}
  \STATE Compute $F(x)$ and $\widehat L_i(x)$. \textbf{If} $F(x)=0$, \textbf{return} $x$.
  \STATE $D\gets C\cap\ball_r(x)$,
     $\Gamma\gets\max\{\Gamma_{\min},8\norm{F(x)}/r\}$.
  \WHILE{true}
    \STATE Form $F_\Gamma$ and compute $L_\Gamma$,
       $m_\Gamma\gets\ceil{\log_2(128L_\Gamma/\Gamma)}$.
    \STATE Split the pairs into $2n$ levels and add the regularizer to one function.
    \STATE Run $m_\Gamma$ passes on fixed $D$, starting from $x$ and reusing each output as the next starting point. Use $\delta=\Gamma/24$ and the known local constants. Let $y$ be the final output.
    \IF{$\norm{y-x}\ge r/4$ or $\Gamma=\Gamma_{\min}$}
      \STATE \textbf{break}
    \ENDIF
    \STATE $\Gamma\gets\max\{\Gamma_{\min},\Gamma/2\}$.
  \ENDWHILE
  \STATE $u\gets\Pi_D(y-F_\Gamma(y)/L_\Gamma)$,
     $n_D\gets L_\Gamma(y-u)-F_\Gamma(y)$.
  \STATE $v\gets F(u)+n_D$, $\lambda\gets1/\Gamma$,
     $x\gets\Pi_C(x-\lambda v)$.
  \STATE $S\gets S+\lambda$, $W\gets W+\lambda u$.
\ENDWHILE
\RETURN $W/S$.
\end{algorithmic}
\end{algorithm}


\subsection{Глобальный зазор и число внешних шагов}

Положим $\lambda=1/\Gamma$, $e_H=u-x+\lambda v$ и
$x^+=\Pi_C(x-\lambda v)$. Из выпуклости, монотонности и
наличия нормали получаем $G(u,w)\le\ip{v}{u-w}$.
Раскрывая квадраты и пользуясь свойством проекции, находим
\begin{align}
 2\lambda G(u,w)
 &\le 2\lambda\ip{v}{u-w} \notag\\
 &\le \norm{x-w}^2-\norm{x^+-w}^2-\norm{u-x}^2+\norm{e_H}^2
 \qquad(w\in C).
 \label{eq:hpe-energy}
\end{align}
При $w=z^\star$ первый член слева неотрицателен.
Каждый длинный шаг уменьшает квадрат расстояния до $z^\star$
хотя бы на $(440/441)(29r/128)^2$. Для относительной
точности $\rho$ берём $\Gamma_{\min}=\rho$ и останавливаемся
при первом $S=\sum_k\lambda_k\ge 1/(2\rho)$. Если краткая
проба принята на нижнем пороге, её вес $1/\rho$ даёт остановку.
Так число принятых шагов составляет $\mathcal O(\mathcal K)$,
а $S\le3/(2\rho)$.

\begin{theorem}[Зазор локализованного HPE]\label{thm:hpe}
При приведённом критерии остановки взвешенный выход
$\bar u=S^{-1}\sum_k\lambda_k u_k$ удовлетворяет
\[
 G(\bar u,w)\le\frac{\norm{x_0-w}^2}{2S}
             \le \rho\norm{x_0-w}^2\qquad(w\in C).
\]
При $\rho=\varepsilon/(2R^2)$ это даёт абсолютную точность
$\varepsilon$ на $C\cap\ball_R(x_0)$.
\end{theorem}

Действительно, в \eqref{eq:hpe-energy} последние два члена в сумме
неположительны по \eqref{eq:hpe-res}. При суммировании промежуточные
квадраты сокращаются. Выпуклость $G(\,\cdot\,,w)$ обеспечивает
оценку для взвешенного среднего. Если уже найдено точное решение,
его зазор неположителен.

\subsection{Покомпонентные оракульные бюджеты}\label{sec:hpe-cost}

Пусть $A_i=L_{0,i}+L_{1,i}\norm{g_i(z^\star)}$.
Константы $A_i$ определяются решением и используются в анализе.
Знак $\OO$ ниже скрывает множители, логарифмические по параметрам
задачи и точности. Для абсолютного зазора $\varepsilon$ на шаре
сравнения получаем
\begin{align}
 N_{\nabla p_i}
 &=
 \OO\!\left(17^{2n-1}\left[
       \mathcal K(1+\sqrt{L_{1,i}R})
       +\sqrt{\frac{\mathcal K A_iR^2}{\varepsilon}}\right]\right),
       \label{eq:hpe-grad}\\
 N_{Q_i}
 &=
 \OO\!\left(17^{2n-1}\left[
       \mathcal K+\frac{M_iR^2}{\varepsilon}\right]\right).
       \label{eq:hpe-op}
\end{align}
Оценки включают отклонённые пробы. Для относительной точности
$\rho$ вместо $R^2/\varepsilon$ используется $1/\rho$.
Разделение касается ведущих членов точности и параметров конкретных
оракулов. Глубина $2n$, радиус $r$ и логарифмы зависят от всей задачи.

\subsubsection{Почему нет экспоненты по расстоянию}

На принятой пробе обозначим $z=z_\Gamma^D$.
Каждый внешний центр удовлетворяет $\norm{x-z^\star}\le R$.
Из регуляризованного ВН на $C$ следует $\norm{z-z^\star}\le R$:
складывая ВН в $z$ и $z^\star$, получаем
$\ip{x-z}{z-z^\star}\ge0$. Для остатка линейной модели
$D_{p_i}(z,z^\star)=p_i(z)-p_i(z^\star)
-\ip{g_i(z^\star)}{z-z^\star}$ имеем
\[
 0\le D_{p_i}(z,z^\star)\le
 \ip{F(z)-F(z^\star)}{z-z^\star}
 \le\Gamma\norm{z-x}\norm{z-z^\star}\le\Gamma rR.
\]
Первое верхнее неравенство получается суммированием
неотрицательных брегмановских остатков и вклада монотонных
операторов. Для второго используем оба ВН.

Положим $d_i=\norm{g_i(z)-g_i(z^\star)}$.
Функция $p_i(\cdot)-\ip{g_i(z^\star)}{\cdot}$ минимальна
в $z^\star$ на всём пространстве. При $d_i>0$ сделаем в анализе
один шаг по её градиенту из $z$ с длиной шага по времени
$1/[e(A_i+L_{1,i}d_i)]$. Длина отрезка, умноженная на
$L_{1,i}$, не больше единицы. Коэффициент квадратичной модели
на нём ограничен $e(A_i+L_{1,i}d_i)$, поэтому
\[
 \frac{d_i^2}{2e(A_i+L_{1,i}d_i)}\le D_{p_i}(z,z^\star).
\]
Здесь нужна гладкость на всём пространстве, поскольку пробная
точка анализа может покинуть $C$. Отсюда
$d_i\le2eL_{1,i}\Gamma rR+\sqrt{2eA_i\Gamma rR}$ и
$L_{1,i}d_i\le A_i+4eL_{1,i}\Gamma R$, используя $L_{1,i}r\le1$.
Вторая оценка получается разделением на случаи
$L_{1,i}d_i\le A_i$ и $L_{1,i}d_i>A_i$. При $d_i=0$ они
выполняются непосредственно. Теперь $\norm{z-x}\le r$,
поэтому рост коэффициента на отрезке от $z$ до $x$ даёт
\begin{equation}\label{eq:hpe-curv}
 \frac{\widehat L_i(x)}{\Gamma}
 \le e^2\bigl(2A_i/\Gamma+4e L_{1,i}R\bigr).
\end{equation}

\subsubsection{Сумма по пробам и шагам}

За один проход с параметром $\gamma$ уровни $(p_i,0)$ и $(0,Q_i)$
по \eqref{eq:pa-cost} дают разные масштабы
$1+\sqrt{\widehat L_i(x)/\gamma}$ и $1+M_i/\gamma$.
Оба умножаются на $\mathcal O(17^{2n-1}m_\gamma)$ за пробу.
Аналитический регуляризатор меняет градиентный масштаб лишь
на константу. Для любой предыдущей пробной $\gamma\ge\Gamma$
имеем $m_\gamma\le m_\Gamma$. Геометрическое уменьшение $\gamma$
даёт $\sum_\gamma \gamma^{-a}=\mathcal O(\Gamma^{-a})$
для $a=1/2,1$, включая возможное последнее усечение.

Число проб за шаг имеет порядок
$1+\log(2+\norm{F(x)}/(r\Gamma))$. Для него и для $m_\Gamma$
можно брать логарифмы общей величины
\[
 \mathcal L=2+nR/r+\sum_i A_i/\Gamma_{\min}
       +\frac{\norm{F(z^\star)}+R\sum_i M_i}{r\Gamma_{\min}} .
\]
Это следует из \eqref{eq:hpe-curv} и
$\norm{F(x)}\le\norm{F(z^\star)}+
 \sum_i(d_i+r\widehat L_i(x)+M_iR)$.
Базовый член на шаг несёт множитель $\log^2\mathcal L$,
члены с $\widehat L_i$ и $M_i$ несут $\log\mathcal L$.
Один запрос $F(x)$ и до двух запросов для поправки
добавляют постоянную работу на шаг.
Наконец, $\sum_k1/\Gamma_k=S$ и
$\sum_k\Gamma_k^{-1/2}\le\sqrt{KS}$, где $K$ --- число
принятых шагов. Подстановка $K=\mathcal O(\mathcal K)$,
$S=\mathcal O(R^2/\varepsilon)$ и \eqref{eq:hpe-curv}
даёт \eqref{eq:hpe-grad}--\eqref{eq:hpe-op}.

\clearpage
\section{Условия линейной сходимости}\label{sec:linear}

В общей монотонной задаче оценка \eqref{eq:hpe-energy} даёт скорость
зазора $\mathcal O(1/S)$ по накопленному весу. Чтобы получить
геометрическое убывание, нужны условия, при которых вычисленная
невязка служит индикатором расстояния до решения.

\subsection{Оценка расстояния через невязку}

Обозначим $T=F+N_C$, $Z=T^{-1}(0)$ и
$d(x)=\dist(x,Z)=\inf_{z\in Z}\norm{x-z}$.
Предположим, что на принятых пробах нижнего порога выполняется
\begin{equation}\label{eq:error-bound}
 d(u)\le\kappa\,\dist(0,T(u)),\qquad \kappa>0.
\end{equation}
Название этого локального свойства --- метрическая
субрегулярность~\cite{leventhal}. Оно позволяет судить по малой
невязке о близости к множеству решений.
Из \eqref{eq:error-bound} следует $d(u)\le\kappa\norm v$
для любого вычисленного $v\in T(u)$.

Если $F$ $\mu$-сильно монотонен на $C$, то
$\ip{v}{u-z^\star}\ge\mu\norm{u-z^\star}^2$ по монотонности
нормального конуса, и можно взять $\kappa=1/\mu$.
Оценка \eqref{eq:error-bound} также покрывает билинейные игры
без сильной монотонности. Например, для $C=\R^{2d}$,
$F(x,y)=(By,-B^\top x)$ и квадратной невырожденной $B$ имеем
$Z=\{0\}$ и $\norm{F(x,y)}\ge\sigma_{\min}(B)\norm{(x,y)}$.
При этом $\ip{F(x,y)}{(x,y)}=0$ для всех $x,y$.
Именно связь остатка с расстоянием здесь обеспечивает сжатие.

\subsection{Сжатие и явный критерий остановки}

Меняем нижний порог на $\Gamma_0=1/\kappa$ и на нём
останавливаемся при $\kappa\norm v\le\eta$.
В алгоритме~\ref{alg:hpe} этот критерий заменяет остановку
по сумме весов. На выход возвращаем принятую точку $u$.
Пусть $\sigma=1/21$.
Из \eqref{eq:hpe-energy}, \eqref{eq:hpe-res} и монотонности
получаем для всех принятых шагов
\begin{equation}\label{eq:linear-energy}
 d(x^+)^2\le d(x)^2-(1-\sigma^2)\norm{u-x}^2,\qquad
 \norm v\le\Gamma(1+\sigma)\norm{u-x}.
\end{equation}
Для первого неравенства подставляем ближайшее к $x$ решение.
Для второго используем треугольник в определении $e_H$.
На нижнем пороге \eqref{eq:error-bound} даёт
$d(x)\le d(u)+\norm{u-x}\le(2+\sigma)\norm{u-x}$.
Следовательно,
\begin{equation}\label{eq:hpe-q}
 d(x^+)^2\le q\,d(x)^2,\qquad
 q=1-\frac{1-\sigma^2}{(2+\sigma)^2}
   =\frac{1409}{1849}.
\end{equation}
Это коэффициент для \emph{квадрата} расстояния.
Каждый шаг с $\Gamma>\Gamma_0$ проходит проверку по длине и потому
уменьшает квадрат расстояния как минимум на
$c=(440/441)(29r/128)^2$. Таких шагов за весь запуск не больше
$B=\lfloor R^2/c\rfloor$.

На пороге те же неравенства дают
$(\kappa\norm v)^2\le(11/10)d(x)^2$.
Положим
\[
  m=\ceil{\log_{1/q}\max\{1,(11/10)R^2/\eta^2\}}.
\]
В префиксе из $B+m+1$ шагов есть проба на нижнем пороге,
до которой уже произошло не менее $m$ сжатий.
На этой пробе $\kappa\norm v\le\eta$ и $d(u)\le\eta$.
Поэтому число внешних шагов имеет порядок
\begin{equation}\label{eq:linear-steps}
 \mathcal O\!\left(1+R^2/r^2+\log(1+R/\eta)\right).
\end{equation}
Если \eqref{eq:error-bound} выполнено только в
$\ball_a(z^\star)\cap C$, достаточно старта в $\ball_{a/3}(z^\star)$.
Расстояние внешних центров до этого решения не растёт, а
\eqref{eq:linear-energy} в виде оценки относительно конкретного
$z^\star$ даёт $\norm{u-x}\le2\norm{x-z^\star}$.
Значит, все следующие принятые точки лежат в $\ball_a(z^\star)$.

\subsection{Бюджеты и случай сильной монотонности}

Положим $\mathcal K_\eta=1+R^2/r^2+\log(1+R/\eta)$.
Во всех принятых шагах $\Gamma\ge1/\kappa$.
Из локальной оценки \eqref{eq:hpe-curv} и подсчёта проб в
разделе~\ref{sec:hpe-cost} получаем
\begin{align}
 N_{\nabla p_i}
 &=\OO\!\left(17^{2n-1}\mathcal K_\eta
               [1+\sqrt{\kappa A_i}+\sqrt{L_{1,i}R}]\right),
               \label{eq:linear-grad}\\
 N_{Q_i}
 &=\OO\!\left(17^{2n-1}\mathcal K_\eta[1+\kappa M_i]\right).
               \label{eq:linear-op}
\end{align}
Число вызовов зависит от точности логарифмически при фиксированных
параметрах задачи. Масштаб оператора остаётся линейным по $\kappa M_i$,
а градиентный масштаб растёт как $\sqrt{\kappa A_i}$.

При $\mu$-сильной монотонности на пороге $\Gamma_0=\mu$ можно
оставить в \eqref{eq:hpe-energy} член
$2\Gamma_0^{-1}\ip{v}{u-z^\star}\ge2\norm{u-z^\star}^2$.
Сумма $2\norm{u-z^\star}^2+(440/441)\norm{u-x}^2$ не меньше
$(440/661)\norm{x-z^\star}^2$ по взвешенному неравенству
Коши--Буняковского. Поэтому коэффициент в \eqref{eq:hpe-q}
улучшается до $221/661$.

В этом случае вычисленная невязка также сертифицирует
абсолютный зазор на всём $C$. По выпуклости, монотонности
и интегрированию $F$ вдоль отрезка $[w,u]$ имеем
\[
 G(u,w)\le \ip{v}{u-w}-\frac\mu2\norm{u-w}^2
          \le\frac{\norm v^2}{2\mu}\qquad(w\in C).
\]
В первом неравенстве $Q(w)$ заменяется интегралом верхней оценки
по отрезку, а сильная монотонность даёт вычитаемый квадратичный
член. Следовательно, достаточно потребовать
$\norm v\le\sqrt{2\mu\varepsilon}$ для абсолютного зазора
$\varepsilon$.

\subsection{Перезапуски пути А при сильной выпуклости}

Есть и прямое следствие для пути А. Если сумма $p$ является
$\mu$-сильно выпуклой на $C$, то по условию ВН
$G(y,z^\star)\ge(\mu/2)\norm{y-z^\star}^2$.
Теорема~\ref{thm:patha} при $\rho=\mu/8$ даёт
\[
 \frac\mu2\norm{z_{\rm out}-z^\star}^2
 \le G(z_{\rm out},z^\star)
 \le\frac\mu8\norm{x_0-z^\star}^2 .
\]
Один запуск уменьшает расстояние до решения вдвое.
Повторение из последнего выхода даёт коэффициент $2^{-j}$
после $j$ запусков. Достаточно сильной выпуклости суммы функций.
Цена каждого запуска по-прежнему определяется общим масштабом
градиента и оператора внутри пары из \eqref{eq:pa-cost}.


\clearpage
\section{Объём формальной проверки}\label{sec:lean}

\vspace{.1em}
\begin{tabular}{@{}p{.23\linewidth}p{\dimexpr.77\linewidth-2\tabcolsep\relax}@{}}
\toprule
\raggedright Результат & Что проверено в Lean \\
\midrule
\raggedright Путь А & Рекурсивный проход и проверка секущих, поиск кривизны и радиуса,
глобальный относительный зазор, бюджеты попыток, оракулов и проекций.
Главная теорема \code{Family.pathA_of_symmetric}. \\
\addlinespace[.25em]
\raggedright Локализованный HPE & Сжатие внутреннего прохода, число перезапусков
и проб регуляризации, допустимость поправки, относительная невязка,
глобальный зазор для взвешенного среднего и граница числа внешних шагов
для сертифицированной последовательности. \\
\addlinespace[.25em]
\raggedright Линейная сходимость & Сжатие HPE при оценке расстояния через невязку,
улучшенный коэффициент при сильной монотонности и конечный критерий
остановки. Для пути А проверен полный запуск, уменьшающий расстояние
вдвое при сильной выпуклости суммы. \\
\bottomrule
\end{tabular}

Формальная проверка HPE пока опирается на условия для последовательности
принятых шагов. Построение этой последовательности из всего алгоритма,
включая формирование $2n$ уровней, и общий подсчёт оракульных вызовов
даны в текстовом доказательстве разделов~\ref{sec:hpe} и~\ref{sec:linear}.
Связь с исходным кодом дана через файлы \code{Restart.lean},
\code{ProximalHPE.lean}, \code{ProximalHPELinear.lean} и
\code{StrongConvexity.lean} в проекте
\href{https://github.com/xangsten/SLIDE}{SLIDE}. Отдельная проверка
ключевых теорем HPE и сильновыпуклого пути А показала только стандартные
аксиомы Lean.

\begin{thebibliography}{9}
\raggedright
\bibitem{borodich}
E.~Borodich, A.~Gasnikov, D.~Kovalev.
\textit{On Linear Convergence in Smooth Convex-Concave Bilinearly-Coupled
Saddle-Point Optimization: Lower Bounds and Optimal Algorithms}.
ICML, PMLR 267, pp.~5045--5100, 2025.
\href{https://proceedings.mlr.press/v267/borodich25a.html}{PMLR}.

\bibitem{gorbunov}
E.~Gorbunov, N.~Tupitsa, S.~Choudhury, A.~Aliev, P.~Richtárik,
S.~Horváth, M.~Takáč.
\textit{Methods for Convex \texorpdfstring{$(L_0,L_1)$}{(L0,L1)}-Smooth
Optimization: Clipping, Acceleration, and Adaptivity}.
ICLR, 2025.
\href{https://arxiv.org/abs/2409.14989}{arXiv:2409.14989}.

\bibitem{monteiro}
R.~D.~C.~Monteiro, B.~F.~Svaiter.
\textit{On the Complexity of the Hybrid Proximal Extragradient Method
for the Iterates and the Ergodic Mean}.
SIAM J.~Optim., 20(6), pp.~2755--2787, 2010.
\href{https://doi.org/10.1137/090753127}{DOI}.

\bibitem{leventhal}
D.~Leventhal.
\textit{Metric Subregularity and the Proximal Point Method}.
J.~Math.~Anal.~Appl., 360(2), pp.~681--688, 2009.
\href{https://arxiv.org/abs/0902.4200}{arXiv:0902.4200}.
\end{thebibliography}
\end{document}
